% Suggested replacement text for Section 4.3 (runtime study).
We benchmark the ODE baseline and the analytic HiPPO-SVGP on the solar stream under two timing protocols. In `kernel_only`, we measure only the construction of the interdomain kernel terms used by the variational objective. In `full_update`, we measure a complete variational update step consisting of ELBO evaluation, backpropagation, and one optimizer step. For the runtime-versus-$T$ study we fix the solar observation set and vary the effective stream length $T \in \{100, 300, 1000, 3000, 7000\}$ by changing the number of HiPPO discretization steps over the same time span. For the runtime-versus-$M$ study we fix $T=1000$ and sweep the basis size $M \in \{8, 16, 32, 64, 128\}$. Across all runs we keep the solar data source, kernel family, random feature count ($R=1000$), hyperparameters, dtype, device, and random seed fixed. Each plotted point reports the median wall-clock runtime over 20 measured repeats after 5 warmup iterations, with CUDA synchronization immediately before and after each timing block.

The measured runtime trends separate clearly by scaling axis. As $T$ increases from $100$ to $7000$ with $M=64$, the analytic method stays nearly flat in wall-clock time: the median `kernel_only` runtime moves from $126.8$ ms to $85.9$ ms and the median `full_update` runtime from $115.9$ ms to $92.6$ ms. In contrast, the ODE baseline grows substantially with stream length, from $5.0$ ms to $281.1$ ms in `kernel_only` mode and from $10.3$ ms to $298.0$ ms in `full_update` mode. The analytic implementation is therefore not the fastest option at very short streams, where its constant overhead dominates, but it becomes clearly preferable once $T$ is large; at $T=7000$ it is about $3.3\times$ faster for kernel construction and $3.2\times$ faster for the full update step.

The $M$ sweep shows a different trade-off. At small basis sizes the analytic method is faster than the ODE baseline (for example, at $M=8$ the median `kernel_only` runtime is $5.2$ ms versus $40.2$ ms, and the median `full_update` runtime is $10.7$ ms versus $42.8$ ms). Around $M=32$ the gap closes, and by $M=64$ and $128$ the analytic method is slower in both timing modes. These measurements support the theoretical claim that the analytic construction removes the dominant dependence on stream length, but they also show that it does not uniformly dominate the ODE baseline as the basis order grows: the cost of the spherical-Bessel and polynomial-order recursion becomes the main bottleneck at large $M$. We therefore restrict our claim to the empirically supported regime where the analytic construction is most beneficial, namely long streams and small-to-moderate basis sizes.

% Figure caption for runtime_vs_T.
\caption{Runtime versus stream length $T$ on the solar benchmark. The two panels separate `kernel_only` timing from the full variational update step. The observation set, kernel family, random feature count, dtype, device, and hyperparameters are fixed across methods; only the HiPPO discretization length $T$ changes. The analytic method has a larger constant overhead at very short streams, but its runtime remains nearly flat as $T$ increases, whereas the ODE baseline grows strongly with $T$. Points show the median runtime over 20 measured repeats after 5 warmup iterations, and the shaded band denotes one empirical standard deviation.}

% Figure caption for runtime_vs_M.
\caption{Runtime versus basis size $M$ on the solar benchmark. The top panel measures kernel construction only and the bottom panel measures one complete ELBO-based update step. Here $T=1000$ is fixed, together with the solar observation set, kernel family, random feature count, dtype, device, and hyperparameters. The analytic method is faster at small basis sizes, reaches a crossover region around $M\approx 32$, and becomes slower than the ODE baseline by $M=64$ and $128$. Points show the median runtime over 20 measured repeats after 5 warmup iterations, with the shaded band indicating one empirical standard deviation.}
